\begingroup
\renewcommand{\thesection}{13A}
\section{Mode and Phase of Registered Events}\label{sec:mode}
\status{definition on imported native orientation data; finite algebra proved and checked here}
\subsection*{Native input and the common domain}
The native $G_{60}$--$C_5$--$C_3$ orientation mechanism supplies three reflection frames and the inverse order-three orientation pair. It identifies the native designation torsor $\{b,ab\}$ equivariantly with that inverse pair, up to one global orientation reversal \citep[Secs.~8--10]{cave2026orientation}. This is the input used here; it is distinct from the abstract $V_4$ triality of Section~\ref{sec:carrier}.

Use cycle notation on three labels and write
\begin{equation}
\Cevt=(0\ 1\ 2),\qquad \Mevt=(1\ 2),\qquad
\Cevt^3=1,\quad \Mevt^2=1.
\tag{M1}\label{eq:modalgenerators}
\end{equation}
The subscript distinguishes $\Cevt$ from the six-axis conference matrix $C$. Let $\mathcal T$ be the three transpositions in $S_3$. The relative closure-frame set is
\[
\mathscr F=\{(r,\rho)\in\mathcal T^2:r\neq\rho\},\qquad
q(r,\rho)=r^{-1}\rho\in\{\Cevt,\Cevt^{-1}\}.
\]
It has six members. The native correspondence supplies the interpretation of $r$ as closing reflection, $\rho$ as residue, and $q$ as the retained relative orientation. Neither $r$ nor $\rho$ alone determines $q$ over the full relabeling orbit. The action is simultaneous conjugation,
\[
g\cdot(r,\rho)=(grg^{-1},g\rho g^{-1}),
\qquad q(g\cdot f)=gq(f)g^{-1}.
\]

\subsection*{Definition and conjugacy}
\begin{definition}[Mode]\label{def:mode}
Mode is the $S_3$ orientation parity of the registered event cycle. On the registered frame orbit, choose one reference presentation $f_*$ and define
\begin{equation}
\mu(g\cdot f_*)=\sgn(g)\in\{+1,-1\}.
\tag{M2}\label{eq:modedefinition}
\end{equation}
The two modes are the two orientation classes; the reference fixes which class is called positive.
\end{definition}
\begin{proposition}[Registered modal-phase algebra]\label{prop:modealgebra}
The conjugation action on $\mathscr F$ is free and transitive. The native three-cycle and a register-reversing reflection satisfy
\begin{equation}
\Mevt\Cevt\Mevt=\Cevt^{-1},\qquad
\langle\Cevt,\Mevt\rangle\cong S_3.
\tag{M3}\label{eq:modalconjugacy}
\end{equation}
Its six frame presentations split into the two three-element sheets corresponding to
\[
\{1,\Cevt,\Cevt^2\},\qquad
\{\Mevt,\Mevt\Cevt,\Mevt\Cevt^2\}.
\]
\end{proposition}
\begin{proof}
Two distinct transpositions have trivial common centralizer in $S_3$: a permutation commuting with both commutes with their generated group and hence belongs to its trivial center. Thus every frame stabilizer is trivial. Since $|\mathscr F|=|S_3|=6$, each orbit is the full frame set. The expression in \eqref{eq:modedefinition} is therefore well defined.

Conjugating $(0\ 1\ 2)$ by $(1\ 2)$ gives $(0\ 2\ 1)$, its inverse. These permutations generate all six elements of $S_3$. The kernel of the sign character is $A_3=\langle\Cevt\rangle$, so it gives precisely the two sheets. Even relabelings preserve $q$ and odd relabelings exchange its inverse values, in agreement with the imported orientation correspondence.
\end{proof}

\subsection*{Phase, presentation, and event identity}
Each framed presentation has a unique expression $\Cevt^p\Mevt^m\cdot f_*$, with $p\in\Z_3$ and $m\in\Z_2$. Phase is $p$ and Mode is $(-1)^m$. Their action is
\begin{equation}
\Cevt:(p,m)\mapsto(p+1,m),\qquad
\Mevt:(p,m)\mapsto(-p,1-m).
\tag{M4}\label{eq:modalcoordinates}
\end{equation}
These are six coordinates with the semidirect-product law $C_3\rtimes C_2$, not commuting independent actions. The two modes carry the same event cycle with opposite orientation.

The event identity retained by this construction is its declared invariant orbit under the presentation action. This is a statement about registered presentation, not an identification of successive completed occurrences. A general $S_3$ action need not have six-point orbits; the six-point assertion above follows from the explicit free frame action. Absolute plus/minus names depend on the reference, while the two-class structure and the relative sign between presentations do not.

The $b/ab$ pair supplies orientation \emph{states}; $\Mevt$ is an \emph{operation} exchanging them. The relative product $q$ has order three even though its orientation torsor has two members. None of these objects is a binary reduction of the integer time register by definition.

\subsection*{Verification scope}
The accompanying modal check enumerates all six group elements, six frames, and 36 action--frame pairs. It checks conjugacy, freeness, the relative decoder, the parity partition, and the coordinate action. The native provenance is the cited orientation theorem; the present replay verifies its finite frame algebra. A map from this frame action to the $r_1$ time register, the Born coefficient space, or an orientation-sensitive transport must specify its domain and covariance separately.
\endgroup
\setcounter{section}{13}
